\documentclass[11pt,letterpaper]{article}
\usepackage[T1]{fontenc}
\usepackage{lmodern}
\usepackage[margin=1.05in]{geometry}
\usepackage{amsmath,amssymb,amsthm,booktabs,microtype,xurl}
\usepackage[hidelinks,pdftitle={Nonhalting and PA-independence for a 196-state Turing machine},
  pdfauthor={James Chang}]{hyperref}

\newtheorem{theorem}{Theorem}[section]
\newtheorem{lemma}[theorem]{Lemma}
\newcommand{\PA}{\mathrm{PA}}
\newcommand{\N}{\mathbb{N}}
\newcommand{\Con}{\operatorname{Con}}
\newcommand{\Prov}{\operatorname{Prov}}
\newcommand{\ax}{\mathsf{AX}}
\newcommand{\mpc}{\mathsf{MP}}
\newcommand{\shift}{\mathsf{S}}
\newcommand{\Ic}{\mathsf{I}}
\newcommand{\Mc}{\mathsf{M}}
\newcommand{\Uc}{\mathsf{U}}
\newcommand{\Ec}{\mathsf{E}}
\newcommand{\Nc}{\mathsf{N}}
\newcommand{\Rc}{\mathsf{R}}
\newcommand{\file}[1]{\texttt{\nolinkurl{#1}}}
\setlength{\emergencystretch}{2em}

\title{Nonhalting and PA-independence\\for a 196-state Turing machine}
\author{James Chang}
\date{}

\begin{document}
\maketitle

\begin{abstract}
We give an explicit two-symbol Turing machine with 196 working states
that searches for a contradiction in ordinary first-order Peano
arithmetic. It runs forever on a blank tape, and Peano arithmetic
proves neither its exact arithmetic halting sentence nor its negation.
The axiom command uses a transposed code for universal quantification
and a guarded equality-vacuity clause. The latter replaces the need
for a primitive transitivity head. The modified calculus, the simulation
of the register program, and both unprovability directions for the
literal transition table are formalized in Lean~4.
\end{abstract}

\section{The result}\label{sec:result}

Write $M_{209}$ for the machine of~\cite{pa209} and $M$ for the
machine of Section~\ref{sec:machine}. Let $H$ be the arithmetic halting sentence
of $M$, built from the rows of $M$ as the halting sentence of
$M_{209}$ is built in~\cite[\S5]{pa209}. Ordinary first-order Peano
arithmetic is denoted by $\PA$.

\begin{theorem}\label{thm:main}
If $\PA$ is consistent, the machine $M$ does not halt from the blank
tape and $\PA\nvdash\neg H$. If $\PA$ is sound for the standard model,
also $\PA\nvdash H$.
\end{theorem}

The register program of $M$ differs from that of $M_{209}$ only in the
body of one command (Lemma~\ref{lem:program}). Sections~\ref{sec:sound}
and~\ref{sec:complete} show that the new command calculus is sound and
that the refutation transport of~\cite{pa209} still reaches it.
Sections~\ref{sec:layout} and~\ref{sec:simulation} deal with the table.

The Lean endpoints use Foundation's standard model of ordinary PA to
supply consistency and soundness in the metatheory. They prove blank-tape
nonhalting and both unprovability directions without additional hypotheses.

\section{The machine}\label{sec:machine}

The tape is infinite in both directions, with alphabet $\{0,1\}$ and
initially all cells zero. The working states are $0,\ldots,195$, state
$0$ is initial, and $196$ terminates the run. A row $(q,b,w,D,q')$
reads $b$ at state $q$, writes $w$, moves one cell in direction $D$,
and enters $q'$. There is one row for every pair in
$\{0,\ldots,195\}\times\{0,1\}$. The 392 rows are given in
\file{data/pa196.tm}, with SHA-256 digest
\begin{quote}\small
\nolinkurl{e8bbccd4cfa325251ca95066db116b5aa766e404b0971c33c02772be1c3617a8}.
\end{quote}
The first row is $(0,0,1,R,1)$, and $(63,1,1,R,196)$ is the only row
that enters the halt label.

Like $M_{209}$, the machine consists of a tape engine that interprets
a register program stored as a binary tree. A configuration of the
register program is a \emph{frame}: a marker, the path of the current
site in the tree, a gap, and a register file in which a value $v$ is
the block $0\,1^{v+1}$~\cite[\S3]{pa209}. The 36 states of the engine
have the same rows in $M$ as in $M_{209}$, after renaming. The root of
the tree is read by state $75$. Of the 196 states, 36 form the engine,
13 route leaf instructions to it, and 147 read the 159 distinct nodes
of the tree.

\section{The register program}\label{sec:steps}

The register program of $M_{209}$ enumerates words over three commands
acting on a store $(t,w,d,h)$ and halts when a Shift sees formula code
$1$~\cite[\S4]{pa209}. Shift and $\mpc$ are kept. The command
$\ax(c,x,y)$ is changed in three steps, and stored differently in a
fourth.

\subsection{Step 1: the code of the quantifier}

Let $p(a,b)=\binom{a+b+1}{2}+b$. In $M_{209}$ the four constructors
are $4p(a,b)+r$ with $r=0,1,2,3$ for implication, membership, the
universal quantifier and equality. Here
\[
 \Ic(a,b)=4p(a,b),\quad \Mc(a,b)=4p(a,b)+1,\quad
 \Uc(a,B)=4p(B,a)+2,\quad \Ec(a,b)=4p(a,b)+3 .
\]
The code of the quantifier exchanges the two pairing arguments. Code $0$ is truth and code
$1=\Mc(0,0)$ is the contradiction, as before. Write
$\Nc(B)=\Ic(B,1)$, $B\wedge C=\Nc(\Ic(B,\Nc(C)))$,
$B\leftrightarrow C$ for $\Ic(B,C)\wedge\Ic(C,B)$, and
$\exists x_a\,B$ for $\Nc(\Uc(a,\Nc(B)))$. In the formulas below an
arrow is $\Ic$.

\subsection{Steps 2 and 3: the heads}

In $M_{209}$ the command pushes the saved top, a generalization, the
restricted propositional head $\rho_8$, seven further heads
$\rho_1,\ldots,\rho_7$, and leaves a guarded gate in $t$. Here
$\ax(c,x,y)$ puts $D=d+x$, $H=h+y$ and $A=c$. It pushes the saved top
$t$, the generalization $\Uc(d,h)$, the head
\[
 \rho_8(D,A,H,d)=\bigl((H\to d)\to D\bigr)\to\bigl((D\to H)\to(A\to H)\bigr),
\]
and then the six formulas of Table~\ref{tab:heads}, in that order.

\begin{table}[ht]
\centering\small
\begin{tabular}{@{}cll@{}}
\toprule
Push & Formula & Head of~\cite{pa209} \\
\midrule
1 & $(\exists x_D\,\Uc(D,H))\to H$ & $\rho_6$ \\
2 & $\Uc(D,\Uc(A,H))\to\Uc(A,\Uc(D,H))$ & $\rho_4$ \\
3 & $\Mc(H,D)\to(\Ec(H,A)\to\Mc(A,D))$ & $\rho_5$, antecedents exchanged \\
4 & $\Ec(H,D)\to(\Mc(A,H)\to\Mc(A,D))$ & $\rho_3$ \\
5 & $\exists x_D\,\Ec(H,D)$ & $\rho_1$, equation mirrored \\
6 & $\Uc(D,A)\to(\Uc(D,A\to H)\to\Uc(D,H))$ & $\rho_2$, antecedents exchanged \\
\bottomrule
\end{tabular}
\caption{The six heads pushed by $\ax$ after the saved top,
$\Uc(d,h)$ and $\rho_8$.}\label{tab:heads}
\end{table}

It leaves in $t$ the guarded gate: $0$ if $A=H$ or $A=D$, and
otherwise $F(H,D,A)$, where for labels $a,b$ and a code $c$
\begin{align*}
 \Delta_c(a,b)&=\Ec(a,b)\to\Uc(c,\Ec(a,b)), &
 V&=\Mc(a,b)\to\Uc(c,\Mc(a,b)),\\
 P&=c\to\Mc(c,b), &
 S&=\Uc\bigl(b,\ (\Mc(a,b)\to P)\wedge(\Mc(c,b)\to c)\bigr),\\
 Q&=\Uc(c,\,S\to\Nc(P)), &
 F(a,b,c)&=\Nc\bigl(\Delta_c(a,b)\to(V\to Q)\bigr).
\end{align*}
Finally it sets $d=h=0$. Nine items are pushed, the saved top
included. In $P$, $S$ and $Q$ the number $c$ occurs both as the label
of the variable $x_c$ and as the code of a formula.

The gate of~\cite{pa209} is $G(a,b,c)=\Nc(V\to Q)$:
membership vacuity $V$ together with pointed separation $\Nc(Q)$,
under the guard $c\ne a,b$. Step~2 removes the seventh head
\[
 \rho_7(a,b,c)=\Ec(c,a)\to\bigl(\Ec(c,b)\to\Uc(\Mc(a,b),\Ec(a,b))\bigr)
\]
and puts the vacuity of equality $\Delta_c(a,b)$ under the same guard.
Step~3 is the rearrangement recorded in the third column of
Table~\ref{tab:heads} and in the parameters of $\rho_8$, whose second
and third arguments are exchanged.

The body of $\ax$ is straight-line code: a sequence of 201 items. An
item is a call of one of the macros of~\cite{pa209} (push the
accumulator; push it and load one of the three registers; pop and form
$4p$; the two parameter reads; the two parts of the guard), a single
increment of the accumulator, or the closing jump. The whole program
has $3{,}812$ sites, compared with $3{,}948$ for $M_{209}$.

\begin{lemma}[Program]\label{lem:program}
Outside the body of $\ax$ the register program $R$ of $M$ agrees with
that of $M_{209}$ after translation of addresses. The body of $\ax$ is a concatenation of blocks, each
identical to a macro block of $M_{209}$ and left only to the next
block. It computes the saved top, $\Uc(d,h)$, $\rho_8(D,A,H,d)$, the
six formulas of Table~\ref{tab:heads} and the gate $F(H,D,A)$, in the
codes of Step~1.
\end{lemma}

\begin{proof}
Flattening the program gives 169 sites outside the body that agree
with $M_{209}$ after translation of addresses. The 201 blocks use the
same macros and internal jump offsets, with the closing jump
redirected to the corresponding successor. Their store effects follow
from the macro contracts. Executing the sequence of calls symbolically
gives nine stack entries and a top that are the stated formulas as
trees over $4p(\cdot,\cdot)+r$, with the pairing arguments in the order
of Step~1. The checked source program and its axiom-command binding
are in \file{PA196/Program.lean} and \file{PA196/Source/AXEffect.lean}.
\end{proof}

\section{Soundness}\label{sec:sound}

\begin{lemma}\label{lem:sound}
If $\PA$ is consistent, $R$ does not halt from the clean store.
\end{lemma}

\begin{proof}
Interpret membership and equality in a model of $\PA$ as
in~\cite[\S4.2]{pa209}: $\Mc(a,b)$ says that bit $x_b$ of $x_a$ is
set. Each formula of Table~\ref{tab:heads} is a head of~\cite{pa209}
up to the order of two antecedents, or is $\exists x_a(x_b=x_a)$, and
is true under every assignment. So is $\rho_8$, a propositional
tautology. When $c\ne a,b$, the formula $G(a,b,c)$ is true as
in~\cite{pa209}, and $\Delta_c(a,b)$ is true because the variable
$x_c$ does not occur in $\Ec(a,b)$; hence $F(a,b,c)$, which is
propositionally $G\wedge\Delta_c$, is true. When the guard fails the
gate is truth. Modus ponens and generalization preserve truth under
every assignment, and code $1$ is false. The invariant, the halt test
and the scheduler are those of~\cite[\S4.2]{pa209}.

Exchanging the pairing arguments in the code of the quantifier changes
the encoder and the decoder together. Both components of a code
$4p(u,v)+r$ are smaller than the code in either order, so the decoding
recursion and the fact that every label in the formula with code $c$
is below $c$ are unaffected. The latter is what makes $x_c$ fresh in
the gate.
\end{proof}

\section{Completeness}\label{sec:complete}

Call a formula \emph{derivable} if it has a proof from instances of
the six heads, of $\rho_8$ with $u\le a$, and of the gate with
$c\ne a,b$, by modus ponens and generalization. All arguments in this
section are fixed finite templates in the labels and codes involved,
with case distinctions on equalities between numbers. They can be
carried out inside $\PA$ for coded proofs, as the corresponding
arguments are in~\cite[\S6.1]{pa209}.

\subsection{What the seventh head was used for}

In the Lean development of~\cite{pa209} the head $\rho_7$ occurs in
one module, \file{ActualCertificateNeighborhoodPrimitive.lean}, which
derives from it reflexivity and symmetry of primitive equality and the
vacuity of $\Ec(a,b)$ under the label $\Mc(a,b)$. These are used in two
places. In \file{ActualCertificateNeighborhoodCongruence.lean},
symmetry turns the two membership rewrites $\rho_3,\rho_5$ into
rewrites in the opposite direction. In
\file{ActualCertificateNeighborhoodWitness.lean}, symmetry and the
vacuity give $\Ec(a,b)\to\Rc(a,b,z)$ for $z\ne a,b$, where
\[
 \Rc(a,b,z)=\Uc\bigl(z,\Mc(a,z)\leftrightarrow\Mc(b,z)\bigr)\wedge
            \Uc\bigl(z,\Mc(z,a)\leftrightarrow\Mc(z,b)\bigr)
\]
is the equality relation of the interpretation, and then
$\exists x_a\,\Rc(a,c,z)$ from $\rho_1$. Everything after that uses
primitive equality only through these three facts:
\begin{description}
\item[(R)] for all labels $a,b,d$, both
  $\Ec(a,b)\to(\Mc(d,a)\leftrightarrow\Mc(d,b))$ and
  $\Ec(a,b)\to(\Mc(a,d)\leftrightarrow\Mc(b,d))$;
\item[(N)] $\Ec(a,b)\to\Rc(a,b,z)$ for $z\ne a,b$;
\item[(W)] an existential witness for equality at every pair of
  labels.
\end{description}
We show that (R), (N) and (W) are derivable. Symmetry and
transitivity of primitive equality are not derived, and are not
needed.

\subsection{Tools}\label{sec:tools}

\begin{enumerate}
\item Classical propositional reasoning with code $1$ as falsity. The
  implicational part comes from $\rho_8$ as in~\cite{pa209}. Head~1
  with $X=\Uc(D,\Nc(\Uc(D,H)))$ reads $\Nc(X)\to H$, and
  $1\to(X\to1)$ is implicational, so $1\to H$ for every formula $H$;
  double negation elimination follows.
\item Lifting: from a derivable implication $B\to C$, the implication
  $\Uc(x,B)\to\Uc(x,C)$, by generalization and head~6.
\item Vacuity. Say that $B$ is vacuous at $r$ if $B\to\Uc(r,B)$ is
  derivable. The gate gives this for the atoms $\Mc(a,b)$ and
  $\Ec(a,b)$ at every $r\ne a,b$. If $B$ is vacuous at $r$, so is
  $\Nc(B)$: from $\Nc(\Uc(r,B))\to\Nc(B)$, lifted at $r$, and head~1
  with $D=r$, one gets $\Nc(\Uc(r,\Nc(B)))\to B$. Vacuity then passes
  to implications by lifting $\Nc(B)\to(B\to C)$ and $C\to(B\to C)$,
  and to $\Uc(x,B)$ with $x\ne r$ by lifting at $x$ and head~2. So a
  formula built from atoms is vacuous at every label that is none of
  its atom indices or bound labels.
\item Existential elimination: if $B$ is vacuous at $x$, then
  $\Uc(x,A\to B)$ and $\exists x\,A$ give $B$. Indeed $\Nc(B)$ gives
  $\Uc(x,\Nc(B))$ and then $\Uc(x,\Nc(A))$.
\end{enumerate}
Tools 1, 2 and 4 are those of~\cite{pa209}. Tool~3 is the vacuity
argument of~\cite{pa209} with one more kind of atom.

\subsection{The gate}

\begin{lemma}\label{lem:extract}
$F(a,b,c)\to G(a,b,c)$ and $F(a,b,c)\to\Delta_c(a,b)$ are derivable.
\end{lemma}

\begin{proof}
Put $B=(V\to Q)$, so that $F=\Nc(\Delta_c\to B)$ and $G=\Nc(B)$. From
$B$ one gets $\Delta_c\to B$, hence $F\to\Nc(B)$. From $\Nc(\Delta_c)$
one gets $\Delta_c\to B$ by explosion, hence
$F\to\Nc(\Nc(\Delta_c))$, and double negation elimination gives
$\Delta_c$.
\end{proof}

So every instance of the old gate is derivable at the same parameters.
In particular, for a label $u$ and a formula $\varphi$ with code
$c\ne u$, the gate at $(u,u,c)$ gives the separation instance
\begin{equation}\label{eq:sep}
 \exists x_c\,(S\wedge P),\qquad
 P=\varphi\to\Mc(c,u),\qquad
 S\to\Uc\bigl(u,\ \Mc(c,u)\to\varphi\bigr).
\end{equation}
The variable $x_u$ is bound in $S$ and free in $P$.

\subsection{Reverse rewriting}

Heads 4 and 3 give the direct rewrites
$\Ec(a,b)\to(\Mc(d,a)\to\Mc(d,b))$ and
$\Ec(a,b)\to(\Mc(a,d)\to\Mc(b,d))$ for all labels. Let $\chi(t)$ be
one of the atoms $\Mc(d,t)$, $\Mc(t,d)$ and $\Mc(t,t)$, where $d$ is a
label. The direct rewrites give
\begin{equation}\label{eq:direct}
 \Ec(v,y)\to(\chi(v)\to\chi(y))
\end{equation}
for all labels $v$ and $y$; for $\Mc(t,t)$ by rewriting one position
after the other.

\begin{lemma}[Specialization of a negative property]\label{lem:spec}
Let $y$, $v$, $c$ be labels with $v\ne y$ and $c\ne y$, and let
$\chi$ be as above with $d\ne y$. Then
\[
 \Uc\bigl(y,\ \Mc(c,y)\to\Nc(\chi(y))\bigr)\ \to\
 \bigl(\Mc(c,v)\to\Nc(\chi(v))\bigr)
\]
is derivable.
\end{lemma}

\begin{proof}
Write $\psi(t)=\Mc(c,t)\to\Nc(\chi(t))$. Head~4 gives
$\Ec(v,y)\to(\Mc(c,v)\to\Mc(c,y))$, and~\eqref{eq:direct} gives
$\Ec(v,y)\to(\chi(v)\to\chi(y))$; here $d=v$ is allowed, because
heads 3 and 4 have no side condition. So
$\Ec(v,y)\to(\psi(y)\to\psi(v))$ is derivable: the label $y$ occurs in
$\psi(y)$ only in an antecedent and under a negation, and both are
reached by direct rewrites. Lift at $y$ to obtain
$\Uc(y,\psi(y))\to\Uc(y,\Ec(v,y)\to\psi(v))$. The indices of
$\psi(v)$ are among $c$, $v$, $d$, all different from $y$, so
$\psi(v)$ is vacuous at $y$. Existential elimination with the instance
$\exists x_y\,\Ec(v,y)$ of head~5 gives $\psi(v)$.
\end{proof}

\begin{lemma}[Reverse rewriting]\label{lem:rev}
For all labels $a,b,d$,
$\Ec(a,b)\to(\Mc(d,b)\to\Mc(d,a))$ and
$\Ec(a,b)\to(\Mc(b,d)\to\Mc(a,d))$ are derivable.
\end{lemma}

\begin{proof}
For $a=b$ both are instances of $B\to(C\to C)$. Let $a\ne b$, and let
$\chi$ be as above with $d\ne a$. We first derive
\begin{equation}\label{eq:chi}
 \Ec(a,b)\to(\chi(b)\to\chi(a)).
\end{equation}
Let $\varphi=\Nc(\chi(a))$. If its code equals $b$, take
$\Nc(\Nc(\Nc(\chi(a))))$ instead, whose code is larger. Let $c$ be the
code of $\varphi$. Then $c$ exceeds $a$ and $d$ and differs from $b$.
Take the separation instance~\eqref{eq:sep} with $u=a$. By
Lemma~\ref{lem:spec} with $y=a$ and $v=b$,
\[
 S\to\bigl(\Mc(c,b)\to\Nc(\chi(b))\bigr),
\]
and head~4 gives $\Ec(a,b)\to(\Mc(c,a)\to\Mc(c,b))$. Hence
$(S\wedge P)\to C$ is derivable, where $C$ is the
formula~\eqref{eq:chi}: if $\Ec(a,b)$ and $\chi(b)$ hold and
$\chi(a)$ fails, then $\varphi$ holds, $P$ gives $\Mc(c,a)$, the
rewrite gives $\Mc(c,b)$, and $S$ gives $\Nc(\chi(b))$. The indices
of $C$ are $a$, $b$ and $d$, all different from $c$, so $C$ is vacuous
at $c$. This is where the vacuity of the equality atom is used.
Existential elimination with~\eqref{eq:sep} gives $C$.

For $d\ne a$, \eqref{eq:chi} with $\chi(t)=\Mc(d,t)$ and with
$\chi(t)=\Mc(t,d)$ are the two statements. For $d=a$ the statements
are $\Ec(a,b)\to(\Mc(a,b)\to\Mc(a,a))$ and
$\Ec(a,b)\to(\Mc(b,a)\to\Mc(a,a))$. Head~3 gives
$\Ec(a,b)\to(\Mc(a,b)\to\Mc(b,b))$, head~4 gives
$\Ec(a,b)\to(\Mc(b,a)\to\Mc(b,b))$, and~\eqref{eq:chi} with
$\chi(t)=\Mc(t,t)$ gives $\Ec(a,b)\to(\Mc(b,b)\to\Mc(a,a))$.
\end{proof}

In the semantics the argument is the familiar one: by separation there
is a set $x_c$ that contains $x_a$ if $x_a$ lacks the property $\chi$
and contains only elements that lack it; $x_b$ belongs to every set
that $x_a$ belongs to; so if $x_b$ has the property, $x_a$ has it. The
calculus has no substitution rule. A quantifier is specialized by a
witness from head~5 and rewriting, and the proof is arranged so that
this needs the direct rewrites only.

\begin{lemma}[The defined equality]\label{lem:nbhd}
For $z\ne a,b$, $\Ec(a,b)\to\Rc(a,b,z)$ is derivable.
\end{lemma}

\begin{proof}
The direct rewrites and Lemma~\ref{lem:rev} with $d=z$ give
$\Ec(a,b)\to(\Mc(a,z)\leftrightarrow\Mc(b,z))$ and
$\Ec(a,b)\to(\Mc(z,a)\leftrightarrow\Mc(z,b))$. Lift both at $z$ and
compose with $\Ec(a,b)\to\Uc(z,\Ec(a,b))$ from the gate.
\end{proof}

Lean derives the reverse rewrites and neighborhood equality in
\file{PA196/Logic/Reverse.lean} and
\file{PA196/Logic/Neighborhood.lean}. The arithmetic certificate
construction supplies the corresponding derivations inside arbitrary
PA models.

\subsection{The transport}

\begin{lemma}[Completeness]\label{lem:complete}
$\PA$ proves that if $\Prov_{\PA}(\ulcorner\bot\urcorner)$ then $R$
halts from the clean store.
\end{lemma}

\begin{proof}
\emph{The calculus.} The direct rewrites and Lemma~\ref{lem:rev} give
(R). Lemma~\ref{lem:nbhd} is (N). Head~5 is (W) with the equation
mirrored. Where~\cite{pa209} specializes a quantifier at $z$ to a
label $c$ it assumes $\Ec(z,c)$ and rewrites atoms in both directions;
with the mirrored witness it assumes $\Ec(c,z)$, and (R) again
supplies both directions. For the witness of the defined equality,
existential monotonicity on $\exists x_a\,\Ec(c,a)$ with (N) gives
$\exists x_a\,\Rc(c,a,z)$, and the symmetry of $\Rc$, which is
derivable from its definition, gives $\exists x_a\,\Rc(a,c,z)$. The
old gate is derivable by Lemma~\ref{lem:extract}, and heads 1, 2, 4
and $\rho_8$ are the old heads $\rho_6$, $\rho_4$, $\rho_3$ and
$\rho_8$. Heads 3 and 6 yield $\rho_5$ and $\rho_2$ by exchanging two
antecedents, which is propositional. So the construction
of~\cite[\S6.1]{pa209}, with the module that used $\rho_7$ replaced by
the lemmas above, turns a coded refutation in $\PA$ into a coded
derivation of code $1$ in the present calculus.

\emph{The codes.} The construction uses of the four codes that
encoder and decoder are inverse to each other, that both components of
a code are smaller than the code, the residues modulo $4$, and, for
$\mpc$, the exact value $\Ic(t,y)=4p(t,y)$ that the program computes.
All four hold for the codes of Step~1. Labels that are chosen as codes
of formulas, such as the code of $\varphi$ in Lemma~\ref{lem:rev} and
the padded property code in the use of separation, are formed with the
new constructor and satisfy the same bounds.

\emph{The command words.} With nine pushed items the identities
of~\cite[\S6.2]{pa209} become, writing $\shift$ for Shift,
\[
 \ax(c,x,y)\,\shift^{9}(t,w,d,h)=(\Uc(d,h),w,0,t),
\]
and the following words implement the three stack operations:
\[
\begin{aligned}
\text{pop}\quad & \ax(0,0,0)\,\shift^{10},\\
\text{push}\quad & \ax(0,0,0)\,\shift^{8}\,\mpc(0),\\
\text{discard}\quad & \mpc(1)\,\ax(1,0,0)\,\shift^{9}\,\mpc(0).
\end{aligned}
\]
Every formula tested by
these Shifts is $0$ or at least $2$, because heads and gate have
residue $0$ and $\Uc(d,h)$ has residue $2$. The head in position $j$
from the top, $1\le j\le7$, is selected by discard, an $\ax$ call,
$j+1$ Shifts and $8-j$ pops, and the gate by discard, an $\ax$ call
and nine pops. After discard $d=h=0$, so $(D,H,A)=(x,y,c)$, and the
instance $\rho(\alpha,\beta,\gamma)$ of a head is obtained by reading
its row of Table~\ref{tab:heads} backwards. For $\rho_8(a,b,c,u)$
with $u\le a$, set the offset to $u$ and call $\ax(b,a-u,c)$. The
generalization word is push, the offset $i$, a Shift,
$\ax(i,0,0)$ and $\shift^{9}$; the modus ponens suffix is
$\shift^{2}\,\mpc(b)\,\ax(b,0,0)\,\shift^{10}$. With these words the
induction on fuel and the scheduler argument
of~\cite[\S6.2--6.3]{pa209} apply as written.
\end{proof}

\section{The layout}\label{sec:layout}

A node of the program tree is read by a state that moves right to the
state of its left or right child according to the path bit. Leaves are
instructions, and equal subtrees are stored once. The 201 items of the
body of $\ax$ are the leaves of a binary tree with 68 distinct
internal nodes.

As in $M_{209}$, two nodes may share one state when one of them occurs
only as a left child and the other only as a right child. On a side
where their children differ, the shared state does not move right. It
moves left, onto the cell that holds the bit by which the node was
entered. A second state re-reads that bit and moves right into a state
whose row for the current cell leads to the correct child. The second
and third states read other nodes of the tree, or are the first state
of a macro, so the device costs nothing. Table~\ref{tab:shared} lists
the six pairs of this kind inside the body. The 68 nodes are therefore
read by 62 states.

\begin{table}[ht]
\centering\small
\begin{tabular}{@{}clll@{}}
\toprule
State & Left-child node & Right-child node & Rows sent left \\
\midrule
138 & $(\mathtt{D}\mid\mathtt{HC})$ & $(\mathtt{PiCDCiiPiC}\mid\mathtt{HC})$ & $0\to168$ \\
140 & $(\mathtt{C}\mid\mathtt{C})$ & $(\mathtt{H}\mid\mathtt{C})$ & $0\to158$ \\
160 & $(\mathtt{i}\mid\mathtt{ADCi})$ & $(\mathtt{i}\mid\mathtt{C})$ & $1\to185$ \\
174 & $(\mathtt{i}\mid\mathtt{ACiiCHDCi})$ & $(\mathtt{DCiiHDCii}\mid\mathtt{ACiiCHDCi})$ & $0\to172$ \\
179 & $(\mathtt{A}\mid\mathtt{ADCi})$ & $(\mathtt{CC}\mid\mathtt{ADCi})$ & $0\to142$ \\
183 & $(\mathtt{C}\mid\mathtt{PiC})$ & $(\mathtt{H}\mid\mathtt{DCi})$ & $0\to158$, $1\to167$ \\
\bottomrule
\end{tabular}
\caption{States of $M$ that read two nodes of the body of $\ax$. A
node is written as its two children, each as a string of items:
\texttt{P} pushes, \texttt{D}, \texttt{H}, \texttt{A} push and load a
register, \texttt{C} pops and forms $4p$, \texttt{i} adds one. The
last column gives each symbol whose row moves left and the state it
enters.}\label{tab:shared}
\end{table}

Outside the body, two states of $M_{209}$ that only served to tell
apart nodes of its old body are gone. Each read a single symbol and
stepped back onto the cell just written. The row that entered it now
does directly what the next state did.

\section{Simulation}\label{sec:simulation}

Number the cells of a frame from the first path cell. In a
\emph{clean} frame the path is followed by zeros up to a single header
$1$, which is followed by the register file. The \emph{gap} is the
position of the header. The deepest site of $R$ has a path of 38
bits. The formal simulation uses the reachable frames whose gap is at
least 48 and whose file has at least ten registers. The clean startup
frame has this gap, and the simulation preserves the lower bound.

\begin{lemma}[One instruction]\label{lem:contract}
For every site of $R$ and each outcome of a conditional decrement,
the machine $M$ started in state $75$ on a valid clean frame of
that site with gap at least 48 reaches, in at least one step, the clean frame of the
successor site with the register effect of the instruction. At the
halt site it enters the halt label.
\end{lemma}

\begin{proof}
A run from a clean frame alternates between walks of the states that
are not engine states, and calls of the engine. The engine rows are
those of $M_{209}$. Its register lemmas therefore hold for all
register contents and all gaps as in~\cite[\S3]{pa209}: filling the
gap, marking a register, the increment, the decrement with both
outcomes, the update of the path, and the return scan. The return scan
ends on the cell after the first $1$ that has seven zeros to its left.
For most sites that cell is the root. For 19 sites one pass ends on a
later path cell, and the walk that follows starts there; $M_{209}$ has
the same feature.

In every run that was computed, the states outside the engine read
only path cells and the two cells after the path, and never a cell to
the left of the root. For every site and outcome the sequence of
states, cells and symbols of the walks was computed for several
register contents and found to be the same. Each contract is therefore
a fixed finite composition of rows and engine calls, and it is
determined by one run.

These runs were carried out for all $3{,}812$ sites, with both
outcomes at each of the $1{,}147$ decrement sites, with varied
register contents and at gaps $40$, $41$, $48$ and $61$. All
$68{,}616$ runs end in the required frame, and every one of the 392
rows is read in at least one of them.
\end{proof}

Lean verifies the positive-length contract at each site for the
reachable gap bound, including both outcomes of a conditional
decrement. The coverage theorem is in \file{PA196/Coverage.lean}.

\begin{lemma}[Startup]\label{lem:startup}
From the blank tape $M$ reaches after exactly $58{,}154$ steps the
clean frame of site $25$ with gap $48$, candidate register $2$, every
other logical register $0$, and $26$ further registers equal to $1$.
\end{lemma}

\begin{proof}
By running the machine. The tape after $58{,}154$ steps equals the
start frame of $M_{209}$ in~\cite[\S2]{pa209} cell for cell, with the
head on the first path cell in state $75$. Lean checks the
segments and composes them into the literal startup trace.
\end{proof}

Site $25$ is the first instruction of the scheduler, and the store
with candidate $2$ and all other registers zero is the clean store
of~\cite{pa209}. The finite startup and its connection to the clean
store are proved in \file{PA196/Startup.lean} and
\file{PA196/Nonhalting.lean}.

\begin{proof}[Proof of Theorem~\ref{thm:main}]
By Lemmas~\ref{lem:contract} and~\ref{lem:startup}, $M$ halts from the
blank tape if and only if $R$ halts from the clean store. If $\PA$ is
consistent, $R$ does not halt by Lemma~\ref{lem:sound}, so $M$ does
not halt and $\neg H$ is true in $\N$.

In an arbitrary model of $\PA$, the refutation transport compiles a
coded PA refutation into a halting command word. The scheduler reaches
that word from candidate $2$. The coded tape simulation prepends the
startup, concatenates the instruction traces, and projects the internal
control coding to the literal rows. The resulting trace witnesses $H$
even when the proof code and trace length are nonstandard. This gives
\[
   \PA\vdash \Prov_{\PA}(\ulcorner\bot\urcorner)\to H,
   \qquad
   \PA\vdash \neg H\to\Con(\PA).
\]
G\"odel's second incompleteness theorem~\cite{goedel} excludes a
$\PA$ proof of $\neg H$. If $\PA$ is sound, it does not prove the
false sentence $H$.
\end{proof}

\section{Lean verification}\label{sec:status}

The checked physical simulation binds all 392 transitions, the startup,
and the contracts at all 3,812 source sites. The finite simulation
checks at gaps 40 and 41 are supplementary. The formal contracts use
the reachable lower bound 48, which suffices for the blank-tape theorem.

The internal arithmetic trace construction is in
\file{PA196/Arithmetic/RefutationHalting.lean}. The exact trace sentence
uses the 196 working states and the distinct halt label 196.
\file{PA196/Independence.lean} derives the internal consistency
implication and two-sided independence. The machine theorem is in
\file{PA196/Nonhalting.lean}. \file{Challenge.lean} states the two
headlines with an explicit PA proof calculus and the literal table.
\file{PA196/Palomar/Solution.lean} proves the matching declarations and
relates this presentation to Foundation and the tape model. The
guarded axiom audits cover both presentations.

The machine table and checked proof sources are distributed together.
See \file{docs/reproduction.md} for finite checks, the Lean build,
and the Comparator invocation.

\begin{thebibliography}{9}

\bibitem{pa209}
J. Chang.
Nonhalting and PA-independence for a 209-state Turing machine.
\href{https://github.com/jamesyc/BusyBeaver/tree/69d98241a721ead20b4f0b17df9062950b40fbfe}{Public PA209 snapshot}.

\bibitem{goedel}
K. G\"odel.
\"Uber formal unentscheidbare S\"atze der Principia Mathematica und
verwandter Systeme I.
\emph{Monatshefte f\"ur Mathematik und Physik} 38 (1931), 173--198.
\url{https://doi.org/10.1007/BF01700692}.

\end{thebibliography}
\end{document}
